\documentclass[14pt]{extarticle}
\usepackage[letterpaper,landscape,left=.4in,right=.4in,top=.62in,bottom=.5in,headheight=18pt,headsep=10pt]{geometry}
\usepackage[T1]{fontenc}
\usepackage{lmodern}
\usepackage{tikz}
\usetikzlibrary{shapes.geometric,arrows.meta}
\usepackage{fancyhdr}
\pagestyle{fancy}
\fancyhf{}
\renewcommand{\headrulewidth}{.4pt}
\renewcommand{\footrulewidth}{0pt}
\fancyfoot[L]{\fontsize{10}{12}\selectfont Bellingham Math Circle / Week 8 / F08-RV-v1}
\fancyfoot[R]{\fontsize{10}{12}\selectfont\thepage}
\setlength{\parindent}{0pt}
\newcommand{\pagehead}[2]{\fancyhead[L]{\fontsize{12}{14}\selectfont Week 8 / #1 / Grades #2}}
\newcommand{\txt}[4]{\node[anchor=north west,inner sep=0pt,text width=#3in,align=left,font=\fontsize{14}{18}\selectfont] at (#1,#2) {#4};}
\newcommand{\smalltxt}[4]{\node[anchor=north west,inner sep=0pt,text width=#3in,align=left,font=\fontsize{11}{14}\selectfont] at (#1,#2) {#4};}
\newcommand{\sheetbegin}{\null\begin{tikzpicture}[remember picture,overlay,x=1in,y=1in]\coordinate (O) at (current page.north west);\begin{scope}[shift={(O)},yscale=-1]}
\newcommand{\sheetend}{\end{scope}\end{tikzpicture}}

% All dimensions are in inches, without scaling the printed page.
\newcommand{\pilecard}[4]{%
  \begin{scope}[shift={(#1,#2)}]
  \node[anchor=north west,font=\fontsize{10}{12}\selectfont,inner sep=0pt] at (0,0) {Start #3};
  \foreach \amount [count=\p] in {#4}{%
    \ifnum\amount>0
      \foreach \n in {1,...,\amount}{\filldraw[fill=gray!25,draw=black,line width=.45pt] ({.36+.48*(\p-1)},{.23+.13*(\n-1)}) circle[radius=.058];}
    \fi
    \node[font=\fontsize{11}{13}\selectfont] at ({.36+.48*(\p-1)},.81) {\amount};
  }
  \node[anchor=north west,inner sep=0pt,font=\fontsize{10}{12}\selectfont] at (1.50,.12) {Winner:};
  \draw[line width=.35pt] (1.50,.62)--(2.05,.62);
  \end{scope}}

% A row has twelve playable squares. The thick boundary is a wall, not square 0.
\newcommand{\coinrow}[4]{%
  \begin{scope}[shift={(#1,#2)}]
  \ifdim#3in>.5in
    \def\numberfont{\fontsize{12}{14}\selectfont}
  \else
    \def\numberfont{\fontsize{8}{9}\selectfont}
  \fi
  \foreach \i in {0,...,12}{\draw[line width=.45pt] ({\i*#3},0)--({\i*#3},#3);}
  \draw[line width=.45pt] (0,0)--({12*#3},0) (0,#3)--({12*#3},#3);
  \draw[line width=2.2pt] (0,0)--(0,#3);
  \foreach \i in {1,...,12}{\node[font=\numberfont,inner sep=0pt] at ({(\i-.5)*#3},{#3+.10}) {\i};}
  \foreach \i in {#4}{\filldraw[fill=gray!65,draw=black,line width=.4pt] ({(\i-.5)*#3},{.5*#3}) circle[radius={.25*#3}];}
  \end{scope}}
\newcommand{\coincard}[5]{%
  \smalltxt{#1}{#2}{1.8}{Start #3}
  \coinrow{#1}{#2+.23}{.23}{#4}
  \smalltxt{#1+3.02}{#2+.20}{1.0}{Winner:}
  \draw[line width=.35pt] (#1+3.02,#2+.61)--(#1+4.22,#2+.61);
}
\newcommand{\activecoins}[1]{%
  \smalltxt{.4}{#1-.30}{4}{Play here}
  \coinrow{.4}{#1}{.85}{}
  \node[anchor=south west,inner sep=0pt,font=\fontsize{9}{10}\selectfont] at (.43,#1-.02) {wall};
}

% Reference and active grids use the same orientation: star at bottom left.
% Token coordinates are distances from the star: x right, y up.
\newcommand{\rookgrid}[5]{%
  \begin{scope}[shift={(#1,#2)}]
  \foreach \i in {0,...,4}{\draw[line width=.5pt] ({\i*#3},0)--({\i*#3},{4*#3}) (0,{\i*#3})--({4*#3},{\i*#3});}
  \pgfmathsetlengthmacro{\starsize}{.43*#3*1in}
  \node[star,star points=5,star point ratio=2.3,draw=black,fill=gray!20,line width=.45pt,minimum size=\starsize,inner sep=0pt] at ({.5*#3},{3.5*#3}) {};
  \ifnum#4>-1
    \filldraw[fill=gray!65,draw=black,line width=.45pt] ({(#4+.5)*#3},{(3.5-#5)*#3}) circle[radius={.24*#3}];
  \fi
  \end{scope}}
\newcommand{\rookpair}[8]{%
  \node[anchor=south,font=\fontsize{10}{12}\selectfont,inner sep=0pt] at ({#1+2*#3},{#2-.045}) {A};
  \node[anchor=south,font=\fontsize{10}{12}\selectfont,inner sep=0pt] at ({#1+6*#3+.16},{#2-.045}) {B};
  \rookgrid{#1}{#2}{#3}{#4}{#5}
  \rookgrid{#1+4*#3+.16}{#2}{#3}{#6}{#7}
}
\newcommand{\rookcard}[7]{%
  \smalltxt{#1}{#2}{2.8}{Start #3}
  \rookpair{#1}{#2+.38}{.31}{#4}{#5}{#6}{#7}{}
  \smalltxt{#1}{#2+1.83}{1}{Winner:}
  \draw[line width=.35pt] (#1+.67,#2+2.05)--(#1+2.64,#2+2.05);
}

\begin{document}
\pagehead{Last counter}{2--5}
\sheetbegin
\txt{.4}{.72}{10.2}{Take turns. Remove one or more counters from one pile. Whoever takes the final counter loses.}
\txt{.4}{1.38}{10.2}{\textbf{Problem 1:} Build these starts with counters. Which player can always win? Find a rule for all two-pile starts with at least one counter.}
\pilecard{.55}{2.18}{1}{1,0}
\pilecard{3.15}{2.18}{2}{1,1}
\pilecard{5.75}{2.18}{3}{1,2}
\pilecard{8.35}{2.18}{4}{2,2}
\pilecard{.55}{3.37}{5}{2,3}
\pilecard{3.15}{3.37}{6}{3,3}
\pilecard{5.75}{3.37}{7}{2,0}
\pilecard{8.35}{3.37}{8}{1,1,1}
\pilecard{.55}{4.56}{9}{1,1,2}
\pilecard{3.15}{4.56}{10}{1,2,2}
\pilecard{5.75}{4.56}{11}{1,2,3}
\pilecard{8.35}{4.56}{12}{2,3,4}
\foreach \x/\n in {.75/1,4.05/2,7.35/3}{%
  \draw[rounded corners=7pt,line width=.7pt] (\x,6.05) rectangle ({\x+2.9},7.45);
  \node[anchor=north west,inner sep=0pt,font=\fontsize{11}{13}\selectfont] at ({\x+.12},6.17) {Pile \n};
}
\smalltxt{.75}{5.73}{3}{Play here}
\smalltxt{.75}{7.62}{1}{Rule:}
\draw[line width=.35pt] (1.3,7.84)--(10.25,7.84);
\sheetend
\newpage

\pagehead{Coin row}{K--3}
\sheetbegin
\txt{.4}{.72}{10.2}{Take turns sliding one coin left through empty squares. Coins cannot jump. A player with no move loses.}
\smalltxt{.4}{1.53}{2.7}{Before}
\smalltxt{4.1}{1.53}{3.2}{Slide the coin from 4 to 2}
\smalltxt{7.8}{1.53}{2.7}{After}
\coinrow{.4}{1.91}{.23}{4,10}
\coinrow{4.1}{1.91}{.23}{4,10}
\draw[-{Stealth[length=4pt]},line width=.8pt] (4.905,1.865)--(4.445,1.865);
\draw[dashed,line width=.5pt] (4.445,2.025) circle[radius=.0575];
\coinrow{7.8}{1.91}{.23}{2,10}
\txt{.4}{2.58}{10.2}{\textbf{Problem 2:} Play these starts. Which player can always win?}
\coincard{.55}{3.24}{1}{1,2}{}
\coincard{5.85}{3.24}{2}{1,4}{}
\coincard{.55}{4.08}{3}{4,5}{}
\coincard{5.85}{4.08}{4}{4,6}{}
\coincard{.55}{4.92}{5}{7,8}{}
\coincard{5.85}{4.92}{6}{7,10}{}
\activecoins{6.25}
\sheetend
\newpage

\pagehead{Coin row}{2--5}
\sheetbegin
\txt{.4}{.72}{10.2}{\textbf{Problem 2 (continued):} Play these three-coin starts. Find a rule for the starts where the second player can always win.}
\coincard{.55}{1.66}{7}{1,2,3}{}
\coincard{5.85}{1.66}{8}{1,4,6}{}
\coincard{.55}{2.59}{9}{2,3,5}{}
\coincard{5.85}{2.59}{10}{2,4,7}{}
\coincard{.55}{3.52}{11}{3,4,7}{}
\coincard{5.85}{3.52}{12}{3,5,9}{}
\coincard{.55}{4.45}{13}{4,7,11}{}
\coincard{5.85}{4.45}{14}{4,7,12}{}
\smalltxt{.55}{5.37}{1}{Rule:}
\draw[line width=.35pt] (1.1,5.62)--(10.1,5.62);
\activecoins{6.25}
\sheetend
\newpage

\pagehead{Two rook boards}{4--5}
\sheetbegin
\txt{.4}{.72}{10.2}{On your turn, choose one board. Slide its token one or more squares left or down. The player who leaves both tokens on their stars wins.}
\smalltxt{.6}{1.47}{3.1}{Before}
\smalltxt{4.05}{1.47}{3.1}{A moves left 2 squares}
\smalltxt{7.5}{1.47}{3.1}{B moves down 2 squares}
\rookpair{.6}{1.98}{.18}{2}{2}{1}{3}{}
\rookpair{4.05}{1.98}{.18}{0}{2}{1}{3}{}
\rookpair{7.5}{1.98}{.18}{0}{2}{1}{1}{}
\txt{.4}{2.96}{10.2}{\textbf{Problem 3:} Play the pairs of starts on the next page. Which pairs let the first player always win? Invent another pair of different starts that lets the second player always win, even though each board alone lets the first player always win.}
\smalltxt{.85}{3.77}{3.4}{A: play here}
\smalltxt{6.2}{3.77}{3.4}{B: play here}
\rookgrid{.85}{4.08}{.85}{-1}{-1}
\rookgrid{6.2}{4.08}{.85}{-1}{-1}
\sheetend
\newpage

\pagehead{Two rook boards}{4--5}
\sheetbegin
\txt{.4}{.72}{10.2}{\textbf{Problem 3 (continued):} Which player can always win from each pair? Use the large boards to play.}
\rookcard{.55}{1.45}{1}{1}{2}{1}{2}
\rookcard{4.0}{1.45}{2}{1}{1}{2}{2}
\rookcard{7.45}{1.45}{3}{1}{0}{2}{3}
\rookcard{.55}{3.90}{4}{1}{3}{2}{2}
\rookcard{4.0}{3.90}{5}{0}{2}{1}{3}
\rookcard{7.45}{3.90}{6}{0}{1}{1}{3}
\smalltxt{.55}{6.27}{2.4}{Your pair:}
\rookpair{3.9}{6.42}{.28}{-1}{-1}{-1}{-1}{}
\sheetend
\end{document}
